\section{Projected vacua and the raw-theta index factor}
\label{sec:sr-source-raw-theta-norm}
This section follows the source normalized isogeny sum, followed
by \srcid{JetTranslateProjectedVacuumLe} and
\srcid{NormRawThetaSeriesIndexBoundLeSqrt}.
It keeps the same Hermitian ambient space, actual Weil action,
frame and algebraic theta reading in every comparison.

\subsection{The normalized isogeny sum on two actual tori}
Use two source Gaussian tori \(Y,Y'\) on the same
\(\mathbb C^N\), with lattices \(L,L'\), unit factors \(k,k'\),
and the actual translation
\[
 (U_aF)(z)=
       e^{\pi H(z,a)-\pi H(a,a)/2}F(z-a).
 \tag{RV1}\label{eq:sr-source-fock-translation}
\]
The factors obey their integral semicharacter law and agree on
\(L\cap L'\).
Let \(R\subset L'\) be a finite set of unique representatives
of \(L'/(L\cap L')\), and set
\[
                     G_R=\sum_{a\in R}k'(a)U_aF
                                      \qquad(F\in Y.\mathrm{Sec}(1)).
 \tag{RV2}\label{eq:sr-source-isogeny-sum}
\]

\begin{srlemma}\label{lem:sr-source-isogeny-sum-norm}
The actual sum \(G_R\) belongs to \(Y'.\mathrm{Sec}(1)\), and
\[
                           L_{2,Y'}(G_R)=|R|\,L_{2,Y}(F).
 \tag{RV3}\label{eq:sr-source-isogeny-sum-norm}
\]
\end{srlemma}
\begin{proof}
The Fock composition law and the semicharacter law say that
\(k'(a)U_a\) is an additive representation for \(a\in L'\).
For \(\gamma\in L'\), write
\(\gamma+a=a'+b\), uniquely with \(a'\in R\) and
\(b\in L\cap L'\).
The operator \(k'(b)U_b=k(b)U_b\) fixes \(F\).
Thus translation by \(\gamma\) permutes the summands in (RV2),
which is the section condition for \(Y'\).
Each summand is entire.

The finite representative set supplies a full common sublattice
\(L_0\subset L\cap L'\): for each real basis vector \(b_i'\)
of \(L'\), two integer multiples have the same representative
modulo \(L\), so a nonzero integer multiple \(n_i b_i'\)
belongs to \(L\); these multiples are still a real basis.
Integrate on its basis parallelepiped.
For \(a,a'\in R\), put
\[
 h_{a,a'}(z)=U_aF(z)\overline{U_{a'}F(z)}e^{-\pi H(z,z)} .
\]
Substitution of (RV1) and the section condition gives exactly
\[
 h_{a,a'}(z+x)
   =e^{2\pi i\,\operatorname{Im}H(x,a-a')}h_{a,a'}(z)
                                             \qquad(x\in L).
 \tag{RV4}\label{eq:sr-source-pairing-density}
\]
For \(x\in L_0\) this multiplier is one, by integrality in \(L'\).
For \(a\ne a'\), their difference is not in \(L\).
Perfection of \(L\) therefore gives some \(x\in L\) for which
the multiplier in (RV4) is not one.
Indeed if every such imaginary pairing were integral,
the resulting integral functional would be represented by
a point of \(L\); nondegeneracy over the real span would then
put \(a-a'\) in \(L\).
The integral of the \(L_0\)-periodic continuous density is
unchanged by any translation of its fundamental domain.
Equation (RV4) consequently forces its off-diagonal integral
to be zero.

On the diagonal,
\(h_{a,a}(z)=|F(z-a)|^2e^{-\pi\|z-a\|^2}\).
Its average over a fundamental domain of \(L_0\) equals
\(L_{2,Y}(F)\), independently of \(a\), by tiling with
fundamental domains of \(L\).
Similarly the average of the density of \(G_R\), periodic
under \(L'\), equals \(L_{2,Y'}(G_R)\).
Expanding the finite square, all off-diagonal averages vanish
and each of the \(|R|\) diagonal terms has the indicated value.
The factors \(k'(a)\) have modulus one.
All integrals are legitimate because the densities are continuous
on bounded fundamental parallelepipeds.
This proves (RV3), by exactly the common-lattice and pairing-density
argument in
\srcid{ContinuousThetaSectionPointwiseNormSq.isogeny_sum_mem_and_norm}.
\end{proof}

\subsection{The normalized tensor jet}
Let \(T\) be the source \(O_E\)-linear positive torus, with
level-one semicharacter determined by the fixed integral \(w_0\)
satisfying \(w_0+\overline{w_0}=1\).
Let \(s\in T.\mathrm{Sec}(1,1)\), \(s\ne0\).
For a positive self-dual Hermitian coefficient lattice
\(\Lambda\subset E^n\), use its actual tensor lattice and
representatives \(R\) modulo the original product period lattice.
Both tensor tori carry the same formula
\[
 \epsilon_{w_0}(\gamma)
 =\exp\{\pi i\,\operatorname{Im}
                         H_C(\gamma,w_0\gamma)\}.
\]
They therefore have the required equality of factors on their
intersection.

The source obtains a real basis of this tensor lattice by
integer commensurability with the product lattice.
Integrality and perfection are transported by its contraction
description: contraction against every coefficient vector in
\(\Lambda\) lies in the original period lattice.
The Hermitian self-duality of \(\Lambda\) supplies a finite dual
family, which reconstructs a tensor from these contractions.
Thus the tensor object used here is the actual
\(\mathrm{GT}\), not merely a collection of unrelated norm bounds.

For \(F=s^{\otimes n}\), the product fundamental domain and
Fubini give \(L_2(F)=L_2(s)^n\).
Lemma~\ref{lem:sr-source-isogeny-sum-norm} gives
\[
                         L_2(G_R)=|R|\,L_2(s)^n .
 \tag{RV5}\label{eq:sr-source-tensor-sum-norm}
\]
Also \(L_2(s)>0\).
If it were zero, its continuous nonnegative norm density would
vanish almost everywhere on one fundamental domain, hence
everywhere there and, by the actual section transformation,
everywhere; this would imply \(s=0\).

Put \(f=[E:\mathbb Q]/2\), \(g=\dim_\mathbb C T=fr\).
Suppose the original period lattice has minimum at least
\(\lambda>0\).
The exact short-vector and contraction argument in
Section~\ref{sec:sr-source-gaussian-sup} gives the tensor
minimum
\(\lambda'=\lambda/(2f^{3/2}n^2\mathrm{rd}_E)\) for \(n>0\);
the zero-row case is handled separately as there.
For a homogeneous polynomial \(P\) of degree \(d\), let
\(\|P\|_1=\sum_\alpha|P_\alpha|\), and define
\[
 {\cal B}(P,d,\lambda)=
 \|P\|_1(\pi e d)^{d/2}\sqrt2
 \left(1+
 {4\sqrt{2gn}\,(2f^{3/2}n^2\mathrm{rd}_E)
                       \over\sqrt\pi\,\lambda}\right)^{gn}.
 \tag{RV6}\label{eq:sr-source-normalized-jet-bound}
\]
The degree-zero factor is \(1\).

Apply the actual Gaussian supremum bound to \(G_R\) on this
tensor torus, then the repeated Cauchy estimate to any translate
\(U_yG_R\). The norm density of \(U_yG_R\) at \(z\) is the
norm density of \(G_R\) at \(z-y\), so its same global bound
applies. For coefficients \(|c_j|\le1\), a finite set \(J\),
and \(|\beta|^2=|R|\), the result is
\[
 { \left|\beta^{-1}\sum_{j\in J}c_j
                         P(\partial)U_{y_j}G_R(0)\right|
                       \over L_2(s)^{n/2}}
                  \le |J|\,{\cal B}(P,d,\lambda).
 \tag{RV7}\label{eq:sr-source-normalized-isogeny-jet}
\]
Here \(\beta\ne0\), since \(R\) contains a representative of zero.
The cancellation is literal:
the square root of (RV5) is
\(|\beta|L_2(s)^{n/2}\).
This is
\srcid{NormalisedThetaValueBound.normalised_theta_value_bound}.
No beta factor is chosen independently of the actual index.

\subsection{The actual projected vacuum at an orthogonal frame}
Keep the source Hermitian datum \(X\), one fixed symplectic
rational basis \(\mathrm{bas}\), and the corresponding genuine
Weil and Heisenberg actions.
Let \(L\subset E^r\) be self-dual for the source alternating
trace form \(EF_X\), \(n\) even, and let \(v\) be the normalized
vacuum for \(\mathrm{tens}(O_E^n,L)\).
For the positive coefficient lattice \(\Lambda\), put
\[
 J_\Lambda=[\Lambda:\Lambda\cap O_E^n],\qquad
                         A_\Lambda=\mathrm{tens}(\Lambda,L).
\]
The source tensor-index formula gives a finite representative set
\(R\subset A_\Lambda\) modulo \(\mathrm{tens}(O_E^n,L)\), with
\[
                  |R|=J_\Lambda^r=:J>0,\qquad
 \mathrm{projA}_{A_\Lambda}v
       ={1\over J}\sum_{a\in R}\epsilon_{w_0}(a)\,\omega(a)v .
 \tag{RV8}\label{eq:sr-source-projector-representatives}
\]
To identify this with the definition as the mean over distinct
orbit functions, use the source vacuum orthogonality:
\(\epsilon(a)\omega(a)v=v\) holds exactly on its original
self-dual lattice. The Heisenberg composition law then gives
equality of two orbit functions exactly when their parameters
differ by that lattice.
Representatives therefore map bijectively to the distinct
orbit functions, including their count \(J\).
This is \srcid{ProjAEqSumReps.projA_eq_sum_reps} and
\srcid{vac_proj_mean}.
The tensor-index formula itself is the source
\srcid{TensorLatticeSelfDualSemicharacterIndex}:
a pseudobasis of either factor makes the tensor relative index
the product of the factor indices; fractional ideal multiplication
does not change each relative index.
It gives
\(\operatorname{relIndex}(\Lambda'\otimes L,\Lambda\otimes L)
 =\operatorname{relIndex}(\Lambda',\Lambda)^r\).
Self-duality gives the required finite nonzero index and vacuum
overlap. We retain the source tensor-vacuum construction used
in this step.

Fix an orthogonal frame \(x\in\mathcal D_X\), and a common
\(t>0\) with \(X.c_\tau=t\) for every \(\tau\).
The source constructs the positive torus \(T\) and
\(\psi:\mathbb C^g\overset{\sim}\longrightarrow V_T\)
from this same \(X,L,x\).
In the actual relabelled coordinates,
\[
 H_T(\psi z,\psi z')=h_x(z,z'),\qquad
 \psi(e_j)=\sqrt t\,e_{\epsilon(j)},\qquad
 \psi(p_v(x,u))\in L_T\ \Longleftrightarrow\ u\in L.
 \tag{RV9}\label{eq:sr-source-self-torus}
\]
For precision, this construction first uses the full lattice
\(L\), an integer commensurability with \(X.\Lambda\), and a
real lattice basis of its scaled period image.
The integral dual of \(L\) represents every lattice functional:
extend it rationally and invert the nondegenerate \(EF_X\)
pairing. Self-duality puts the representing vector back in \(L\).
The torus factor is the sign of
\(\operatorname{Im}H_T(\gamma,w_0\gamma)\);
the integer cross-term formula and \(w_0+\overline{w_0}=1\)
give its semicharacter law.
Finally \(\psi\) is the explicit coordinate permutation followed
by multiplication by \(\sqrt t\).
Uniqueness of the Hermitian form with the given imaginary
pairing on a full period lattice identifies it with \(h_x\).
These are the steps in
\srcid{ExistsTorusOfSelfD.exists_torus_of_selfD}.
The equality of the period coordinates with the source \(h_x\)
uses the specified \srcid{HermitianFormAtFrameLaws}.

Let the actual normalized theta function satisfy
\[
 \theta\in\mathrm{IsTheta1}(X,x,L),\qquad
 \mathrm{MeanOne}(\theta),\qquad
       \prod_b\theta(z_b)=\kappa\,\mathrm{thA}(v,x,z).
 \tag{RV10}\label{eq:sr-source-theta-vacuum-tie}
\]
Under (RV9), \(s(v_1)=\theta(\psi^{-1}v_1)\) is a level-one
section with \(L_2(s)=1\).
Indeed substitution in its Fock transformation gives (GS1),
and the normalized integral is unchanged by the fixed
real linear coordinate change: its Haar scalar cancels between
the integral and volume. Any two fundamental domains give the
same integral of the periodic density.
Thus \(\mathrm{MeanOne}\) is exactly this \(L_2\)-normalization;
it is not an unrelated positive scalar.

Write \(\Phi(y)_b=\psi(p_v(x,y_b))\).
The source proves that \(-\Phi(R)\) is the representative set
for the tensor torus above,
\(\epsilon_{w_0}(-\Phi(a))=\epsilon_{w_0}(a)\),
and Fock translation becomes \(U_{-\Phi(y)}\).
Consequently the theta function for
\(\omega(y)\mathrm{projA}_{A_\Lambda}v\) is
\[
                  \kappa^{-1}J^{-1}\,G_R\circ\psi^{\,n},
\]
with the additional common translation \(U_{-\Phi(y)}\).
If \(\kappa=0\), the target bound is immediate.
Otherwise the displayed \(\kappa^{-1}\) cancels its exterior
\(\kappa\).

Each actual theta-polynomial letter becomes
\(\sqrt t\) times its Euclidean coordinate direction.
The relabelling of letters is injective, so it preserves the
coefficient \(L^1\)-sum.
Homogeneity gives the derivative factor \(t^{d/2}\).
Apply (RV7) with one translation and \(\beta=\sqrt J\).
Since (RV8) divides by \(J\), the exact result is
\[
 |\kappa\,\mathrm{thP}(\omega(y)\mathrm{projA}_{A_\Lambda}v,P,x)|
           \sqrt J
 \le \|P\|_1(\pi e t d)^{d/2}\sqrt2
 \left(1+
 {4\sqrt{2gn}\,(2f^{3/2}n^2\mathrm{rd}_E)
                       \over\sqrt\pi\,\lambda}\right)^{gn}.
 \tag{RV11}\label{eq:sr-source-projected-vacuum-jet}
\]
Here the minimum assumption is exactly
\(\lambda^2\le h_x(p_v(x,u),p_v(x,u)).\operatorname{re}\)
for every \(0\ne u\in L\).
The source \srcid{hx_ortho} supplies this from the same frame's
\(\lambda^2\le\mathrm{firstMinimumSq}(x)\).
Equation (RV11) is
\srcid{JetTranslateProjectedVacuumLe.jet_projA_le}.

\subsection{The exact raw-theta value}
Now use the original \(E,H,D_w,l,X_o,c\), good embedding
\(\sigma:c.F\to\mathbb C\), actual \(\mathrm{IsUnf}\) with
\(r=p+q\ge1\), and actual \(\mathrm{IsNab}\) connection.
The row dictionary in the preceding construction has
\(X.p_0=D_w.q\) and the other block size \(q_0=D_w.p\), exactly
as in \srcid{ExistsAnalyticDataUnitFrameFramedPoint.exists_setup}.
The polynomial directions, signatures and \(\psi\) use that
specified dictionary, including this ordering.

Let \(n\) be even, \(a\) dominant, and let
\(\Lambda=\mathrm{La}\) be the actual positive coefficient
lattice with \(\mathrm{Unimod}\),
\(\mathrm{TrivAt}(l.M)\) and \(\mathrm{LevOK}(N,l.M)\).
Let \(P\in\mathrm{harmL}(c,\Lambda,a)\) have total degree
\(d=\mathrm{DF}(D_w,a)\).
Let \(f\in\mathrm{sec}(X_o,c,a,n)\), and suppose its values on
every source test ring and every frame satisfy the actual
\(\mathrm{thLaValAt}\) formula for these same
\(\Lambda,N,P,\mathrm{nab}\) and coefficient function \(w\).
Assume \(|\sigma(w(q))|\le1\) for every label \(q\).

The source first extends the equality from test rings to the
complex frame: actual \(\mathrm{secBaseChange}\) supplies a
test frame mapping to it; uniqueness and naturality of the
actual short theta value give the same value after pullback.
This is \srcid{thLaValAt_evR}.
It uses all test rings in the premise, not equality at one
preferred complex point.

At an orthogonal frame \(x\), the supplied analytic fibre,
marked level data, theta generator and normalized vacuum give
the exact complex value formula. If
\[
 m=\mathrm{mIdx}(\Lambda,N,r),\quad
 z=\mathrm{zCard}(\Lambda,N,r),\quad
 Q=|\mathrm{Fin}(n)\to\mathrm{LevQ}(D_w,l)|,\quad
 J=J_\Lambda^r,
\]
its scalar is
\[
 {z\over m}\ {1\over Q}\,
 \kappa\,\mathrm{thP}
       \left(\sum_q\sigma(w(q))\,
                  \omega(Y_q)\mathrm{projA}_{A_\Lambda}v,\ P,\ x\right),
              \qquad z=mJ .
 \tag{RV12}\label{eq:sr-source-raw-reading-scalar}
\]
The exact relation \(z=mJ\) follows by multiplying the finite
lattice indices through the intersection with \(O_E^n\);
it is \srcid{zCard_eq}, with \srcid{index_Ln}.
No normalizing \(\mathrm{betaH}^{-1}\) is inserted in this raw
formula.

Put \(B_0\) equal to the right side of (RV11).
Linearity of the actual theta derivative and the triangle
inequality give a bound \(B_0/\sqrt J\) for the normalized
\(Q\)-mean in (RV12).
Combining this with \(z=mJ\) gives
\[
 \left|\operatorname{ev}_\sigma(f;x)\right|
 \le \sqrt{J_\Lambda^r}\,
       \left(\sum_\alpha|\sigma(P_\alpha)|\right)
       (\pi e t d)^{d/2}\sqrt2
 \left(1+
 {4\sqrt{2frn}\,(2f^{3/2}n^2\mathrm{rd}_E)
                       \over\sqrt\pi\,\lambda}\right)^{frn}.
 \tag{RV13}\label{eq:sr-source-raw-value-bound}
\]
Algebraically \(|m^{-1}z|\le J\) even in the source's
zero-denominator convention; the actual finite-index case has
\(m>0\).
Likewise an empty-label mean is zero by that convention.
The actual label group is finite and nonempty.
This proves exactly
\srcid{NormRawThetaSeriesIndexBoundLeSqrt.norm_evR_le}
under the stated actual analytic-fibre/reading suppliers.

\paragraph{Frozen-source ledger and retained geometric comparisons.}
All sources use \srcid{alpoge-source-7cbf7a0c}.
In
\srcid{src/Development/AbelianVarietyTheory/Proofs/CuspFormEvaluations/},
the files are
\begin{itemize}
\item \srcid{ContinuousThetaSectionPointwiseNormSq.lean},
\srcid{9f7ab3413f449c9b93dac642ca2ab0845057ee85};
\item \srcid{NormalisedThetaValueBound.lean},
\srcid{689a12045c213d4d11e73f6660d5ee8a8739f684};
\item \srcid{ProjAEqSumReps.lean},
\srcid{818dccb76df27daa5ee0a1c63a15a27700ebfaf4};
\item \srcid{ExistsTorusOfSelfD.lean},
\srcid{859af98c532c58c34785cd7f6e382a0f49944c71};
\item \srcid{L2OneEqIntegralFundamentalDomainDivVolume.lean},
\srcid{34d7c34bb82c87d1f3b0a8a9ea0771ffa8c0d23e};
\item \srcid{JetTranslateProjectedVacuumLe.lean},
\srcid{4a358a4e34143451dab2d499357ed33d1e51249d};
\item \srcid{ExistsAnalyticDataUnitFrameFramedPoint.lean},
\srcid{903ac3995f8acf7789604f184c54eb7fc7887e45}.
\end{itemize}
The tensor lattice and relative-index proof is
\srcid{src/Development/AbelianVarietyTheory/Proofs/SiegelThetaTransformation/TensorLatticeSelfDualSemicharacterIndex.lean},
blob \srcid{4894f884d42f8c17e88ea2604bcf2fc40817b20d}.
The final raw reading proof is
\srcid{src/Development/AbelianVarietyTheory/Proofs/DivisorProofs/NormRawThetaSeriesIndexBoundLeSqrt.lean},
blob \srcid{d1817b6beba51a259498dc4ae71cb669cf67c874}.

The remaining construction-specific suppliers at (RV10)--(RV12)
are exactly the actual
\srcid{AnalyticFibreUnitFrameUniformisation},
\srcid{FramedPointGeneratorIsThetaLevelOne},
\srcid{ProdLevelOneThetaEqVacuumTheta},
\srcid{DualHalfLevelMarkingExistsAndMatch},
and
\srcid{ShortThetaValueComplexPointEqClassicalTheta},
together with actual \srcid{IsNab}, the symplectic row dictionary,
and the source tensor-vacuum/Heisenberg laws.
The preceding arguments prove the Gaussian and index bounds with
these compatible data; they do not construct separate witnesses
and identify them merely because they have the same names.
